\section{Simultaneous norm equations}\label{sec:norm-tori}

Prescribing a determinant inside a unitary torus is a norm equation.
Later we shall prescribe two norms of the same element: one to a
degree-$m$ field and one to a quadratic field.  This section proves
that the resulting equations satisfy the Hasse principle and weak
approximation when the two fields have the indicated symmetric Galois
group.  The finite-lattice argument is that of
\cite[Lemmas~2.7--2.9]{MPNF}; we reproduce it and supply the
global-function-field duality input.

For a finite separable extension $E/K$, write
$\Res^1_{E/K}\bbG_m$ for the kernel of its norm map.

\begin{proposition}[Simultaneous norms]\label{tor:norm-approximation}
Let $k$ be a global function field, let $J/k$ be a separable quadratic
extension, and let $F/k$ be separable of degree $m\geq3$.
Suppose that the Galois group of the joint normal closure of $F$ and
$J$ is $S_m\times C_2$ under its natural action on the $m$
embeddings of $F$ and its restriction to $J$.
Put $E=FJ$.  Then
\begin{equation}\label{tor:norm-kernel}
 T=\ker\left[
 \Res_{F/k}\bigl(\Res^1_{E/F}\bbG_m\bigr)
 \xrightarrow{\,N_{E/J}\,}\Res^1_{J/k}\bbG_m
 \right]
\end{equation}
is a torus.  Every torsor under $T$ that has a point over every
completion of $k$ has a $k$-point, and its $k$-points are dense in
the product of its points over any finite set of completions.

In particular, let $h\in F^\times$ and $d\in J^\times$ satisfy
$N_{F/k}(h)=N_{J/k}(d)$.  If the equations
\begin{equation}\label{tor:two-norms}
                 N_{E/F}(a)=h,\qquad N_{E/J}(a)=d
\end{equation}
have a solution in $(E\otimes_k k_v)^\times$ for every place $v$,
they have a solution in $E^\times$ approximating any prescribed
local solutions at finitely many places.
\end{proposition}

We first compute the character lattice and then express its
arithmetic obstruction in terms of a finite group.
A \emph{$k$-lattice} is a finitely generated free abelian
group with a continuous action of the absolute Galois group of $k$;
its action factors through a finite quotient.
Let $P$ be the joint normal closure in the proposition and let $c$
generate the second factor of $\Gamma=\Gal(P/k)=S_m\times C_2$.
The character map dual to the norm in Equation~\eqref{tor:norm-kernel}
is the diagonal inclusion
\[
 \bbZ_{\mathrm{sign}}\longrightarrow
       \bbZ^m\otimes\mathrm{sign}_{C_2},\qquad
                       1\longmapsto\mathbf1=(1,\ldots,1).
\]
Here $S_m$ permutes the $m$ coordinates and $c$ acts by $-1$.
The inclusion is primitive, so its cokernel is a lattice.
Consequently the norm is surjective with connected kernel $T$,
and
\begin{equation}\label{tor:lattice}
       X^*(T)=M=(\bbZ^m/\bbZ\mathbf1)
                          \otimes\mathrm{sign}_{C_2}.
\end{equation}

For any $k$-lattice $M$, set
\[
 \Sha^2(k,M)=\ker\left[H^2(k,M)\longrightarrow
                              \prod_v H^2(k_v,M)\right],
\]
and let $\Sha^2_\omega(k,M)$ denote the subgroup of classes whose
localizations vanish at all but finitely many places.  All Galois
cohomology here uses continuous cochains and discrete coefficients.

\begin{lemma}[Detection on cyclic subgroups]\label{tor:cyclic-detection}
Suppose that a finite Galois extension $P/k$ splits the $k$-lattice
$M$, and put $\Gamma=\Gal(P/k)$.  Inflation identifies
$\Sha^2_\omega(k,M)$ with
\begin{equation}\label{tor:cyclic-kernel}
 \ker\left[H^2(\Gamma,M)\longrightarrow
                    \prod_{C\subseteq\Gamma\ \mathrm{cyclic}}H^2(C,M)
       \right].
\end{equation}
If this group vanishes and $T$ is the torus with character lattice
$M$, every everywhere locally soluble $T$-torsor has a rational
point and satisfies weak approximation.
\end{lemma}

\begin{proof}
The action on $M$ over $P$ is trivial.  A continuous homomorphism
from a profinite group to a torsion-free discrete group has finite,
hence trivial, image.  Thus $H^1(P,M)=0$; and, since rational
coefficients have no positive-degree cohomology, the sequence
$0\to M\to M\otimes\bbQ\to M\otimes(\bbQ/\bbZ)\to0$ gives
\[
 H^2(P,M)=H^1\bigl(P,M\otimes(\bbQ/\bbZ)\bigr).
\]
An element on the right is a continuous character with finite image.
Such a character cannot be locally zero at almost every place
unless it is zero, by Chebotarev applied to the finite extension
through which it factors.  Therefore a class in
$\Sha^2_\omega(k,M)$ restricts to zero over $P$.  Inflation--restriction,
together with $H^1(P,M)=0$, then gives a unique preimage in
$H^2(\Gamma,M)$.

At a place $v$ unramified in $P$, the decomposition group is cyclic.
The local inflation map from its degree-two cohomology into
$H^2(k_v,M)$ is injective, again because degree-one cohomology of
the lattice over the splitting field vanishes.  Chebotarev supplies
infinitely many places with each possible Frobenius conjugacy class.
Consequently, a class locally zero almost everywhere restricts to
zero on every cyclic subgroup of $\Gamma$.  Conversely, vanishing
on all cyclic subgroups gives local vanishing at every unramified
place.  This proves Equation~\eqref{tor:cyclic-kernel}.

For the arithmetic conclusion, global torus duality over function
fields gives a perfect pairing of finite groups
\[
                  \Sha^1(k,T)\times\Sha^2(k,M)
                                  \longrightarrow\bbQ/\bbZ
\]
\cite[Theorem~5.2]{DemarcheHarariDuality}.
Here $\Sha^1(k,T)$ is the group of everywhere locally trivial
$T$-torsors.  The same theory gives weak approximation for $T$
when $\Sha^2_\omega(k,M)=0$; explicitly, apply
\cite[Corollary~4.5]{DemarcheHarariHomogeneous} to the reductive
group $T$.  In the notation of these references the associated
complex is $[0\to T]$, and its dual is $[M\to0]$, with $M$ in
degree $-1$, so their $\Sha^1_\omega$ of the dual is precisely
$\Sha^2_\omega(k,M)$.
Since $\Sha^2(k,M)\subseteq\Sha^2_\omega(k,M)$, the vanishing in
the statement gives the Hasse principle as well.  Translation by
a rational point identifies any soluble torsor with $T$ and
transfers weak approximation to it.  These duality results include
the characteristic-primary part; no restriction on the
characteristic of $k$ is needed.
\end{proof}

The arithmetic problem has now reduced to showing that
Equation~\eqref{tor:cyclic-kernel} vanishes for the specific lattice
in Equation~\eqref{tor:lattice}.
We use the interpretation of $H^2(\Gamma,M)$ as extensions of
$\Gamma$ by the abelian group $M$ inducing the specified action.
Restriction to a subgroup is zero exactly when the extension
splits over that subgroup.  A central element acting by $-1$
makes these extensions particularly concrete.

\begin{lemma}[A central sign element]\label{tor:sign-element}
Let $\Gamma$ be a finite group acting on a lattice $M$, and suppose
that $c\in Z(\Gamma)$ has order two and acts as $-1$ on $M$.
Every extension
\[
 1\longrightarrow M\longrightarrow\mathcal E
                       \longrightarrow\Gamma\longrightarrow1
\]
defines a class
$\beta\in H^1(\Gamma/\langle c\rangle,M/2M)$ that is zero
if and only if the extension splits.
If the extension splits over an involution $r\in\Gamma$, then
the value at the image of $r$ of every cocycle representing
$\beta$ lifts to an element $d_r\in M$ satisfying
$(1+r)d_r=0$.
\end{lemma}

\begin{proof}
Write $M$ additively, with its translations placed to the left
of lifts in $\mathcal E$.  A lift $z$ of $c$ is an involution:
$z^2$ belongs to $M$ and commutes with $z$, hence belongs to
$M^c=0$.  Choose lifts $s_g$ of $g\in\Gamma$, with $s_1=1$
and $s_c=z$, and define $f(g,h),d_g\in M$ by
\[
 s_gs_h=f(g,h)s_{gh},\qquad
                   zs_gz^{-1}=d_gs_g.
\]
Conjugating the first equality by $z$ in two ways gives
\begin{equation}\label{tor:sign-cocycle}
                    d_{gh}=d_g+g d_h+2f(g,h).
\end{equation}
Thus $g\mapsto\overline d_g$ is a cocycle with values in $M/2M$.
Its value at $c$ is zero, and $c$ acts trivially on $M/2M$, so
it descends to $\Gamma/\langle c\rangle$.
Replacing $s_g$ by $a_gs_g$ changes $d_g$ to $d_g-2a_g$;
replacing $z$ by $az$ changes it to $d_g+(1-g)a$.
The cohomology class $\beta$ is therefore independent of these
choices.

If $\beta=0$, change $z$ so that every $d_g$ is even, and then
replace $s_g$ by $(d_g/2)s_g$.  These lifts all commute with $z$.
The centralizer of $z$ in $\mathcal E$ maps onto $\Gamma$ and
has kernel $M^c=0$; it is therefore a complement to $M$ and gives
a splitting.  Conversely, a splitting supplies mutually compatible
lifts with $d_g=0$.

Finally, suppose that $r$ has an involutive lift $s_r$.
Conjugating $s_r^2=1$ by $z$ shows that
$(1+r)d_r=0$.  Replacing the cocycle by a cohomologous one changes
this value modulo $2M$ by $(1-r)a$ for some $a\in M$.
The corresponding lift $d_r+(1-r)a$ remains annihilated by $1+r$,
since $r^2=1$.  This proves the last assertion.
\end{proof}

\begin{lemma}[The symmetric permutation lattice]\label{tor:symmetric-lattice}
Let $m\geq3$, let $\Gamma=S_m\times\langle c\rangle$ with
$c^2=1$, and let
\[
 M=(\bbZ^m/\bbZ\mathbf1)\otimes\mathrm{sign}_{C_2},
 \qquad \mathbf1=(1,\ldots,1).
\]
Thus $S_m$ permutes coordinates and $c$ acts by $-1$.
An extension of $\Gamma$ by $M$ that splits over every cyclic
subgroup splits over $\Gamma$.
\end{lemma}

\begin{proof}
Fix an extension with the stated cyclic splittings.  By
Lemma~\ref{tor:sign-element}, it suffices to show that its class
\[
          \beta\in H^1(S_m,Q_m),\qquad
          Q_m=\bbF_2^m/\bbF_2\mathbf1=M/2M,
\]
is zero.  The exact sequence of $S_m$-modules
\[
 0\longrightarrow\bbF_2\mathbf1\longrightarrow
 V_m=\bbF_2^m\longrightarrow Q_m\longrightarrow0
\]
has connecting homomorphism
\[
        \partial:H^1(S_m,Q_m)\longrightarrow H^2(S_m,\bbF_2).
\]
Shapiro's lemma identifies $H^i(S_m,V_m)$ with
$H^i(S_{m-1},\bbF_2)$, where $S_{m-1}$ fixes one coordinate.
Under this identification, the map induced by the diagonal
inclusion $\bbF_2\to V_m$ is restriction to $S_{m-1}$.
In degree one it is an isomorphism: for both $S_m$ and $S_{m-1}$,
including $S_{m-1}=S_2$, homomorphisms to $\bbF_2$ are generated
by the permutation sign.  The cohomology exact sequence therefore
shows that $\partial$ is injective and that its image is the
kernel of restriction to $H^2(S_{m-1},\bbF_2)$.

Let
\[
 1\longrightarrow\langle\zeta\rangle\longrightarrow
       \widehat S_m\longrightarrow S_m\longrightarrow1,
 \qquad \zeta^2=1,
\]
be the central extension represented by $\partial\beta$.
It splits over each point stabilizer, because these are conjugate.
Every transposition therefore has involutive lifts; multiplying
such a lift by $\zeta$ preserves its square.  If $m\geq5$, any
two disjoint transpositions fix a common point.  Their lifts
commute, since their commutator is unaffected by changing a lift
by $\zeta$ and is trivial in a splitting over that stabilizer.
For $m=3$ there is no disjoint pair of Coxeter generators.

For these degrees choose involutive lifts $t_i$ of the adjacent
transpositions $(i,i+1)$, $1\leq i<m$.  Write
\[
                  (t_it_{i+1})^3=\zeta^{\epsilon_i},
                  \qquad \epsilon_i\in\bbF_2.
\]
Replace $t_i$ by $\zeta^{a_i}t_i$, where $a_1=0$ and
$a_{i+1}=a_i+\epsilon_i$.  This makes all adjacent-product cubes
equal to $1$ without changing squares or commutation of distant
generators.  The Coxeter presentation
\[
 s_i^2=1,\qquad (s_is_{i+1})^3=1,\qquad
                   s_is_j=s_js_i\quad (|i-j|>1)
\]
now gives a section $S_m\to\widehat S_m$.  Hence
$\partial\beta=0$, and injectivity gives $\beta=0$.

The degree $m=4$ requires the integral information retained in
Lemma~\ref{tor:sign-element}.  The transpositions
$\tau=(12)$ and $\nu=(34)$ have no common fixed point, so their
lifts might fail to commute.  Choose the cocycle
$g\mapsto\overline d_g$ obtained from the original extension by
the lattice, and choose representatives $b(g)\in\bbF_2^4$ of its
values in $Q_4$, with $b(1)=0$.  The factor set for
$\partial\beta$ is given by
\[
       \epsilon(g,h)\mathbf1=b(g)+g b(h)-b(gh).
\]
As $\tau\nu=\nu\tau$, noncommuting lifts of $\tau$ and $\nu$
would force
\begin{equation}\label{tor:four-obstruction}
          (1+\nu)b(\tau)+(1+\tau)b(\nu)=\mathbf1.
\end{equation}

The original lattice extension splits over $\langle\tau\rangle$.
Lemma~\ref{tor:sign-element} consequently gives a lift $d\in M$
of $\overline d_\tau$ with $(1+\tau)d=0$.
Choose a representative $\widetilde d\in\bbZ^4$.  For some
$a\in\bbZ$,
\[
       \widetilde d+\tau\widetilde d=a\mathbf1,
       \qquad 2\widetilde d_3=a=2\widetilde d_4.
\]
It follows that $\widetilde d_3=\widetilde d_4$.  Thus every
representative of $\overline d_\tau$ in $\bbF_2^4$, in particular
$b(\tau)$, has equal third and fourth coordinates.  The third and
fourth coordinates of $(1+\nu)b(\tau)$ are zero.  Those of
$(1+\tau)b(\nu)$ are also zero because $\tau$ fixes these
coordinates.  This contradicts Equation~\eqref{tor:four-obstruction}.
The lifts of $\tau$ and $\nu$ commute.  These form the only disjoint
pair among the adjacent transpositions in $S_4$, so the same adjustment of adjacent
products proves that $\widehat S_4$ splits.

We have proved $\beta=0$ in every degree $m\geq3$.
Lemma~\ref{tor:sign-element} splits the original extension.
\end{proof}

\begin{proof}[Proof of Proposition~\ref{tor:norm-approximation}]
We computed the character lattice of $T$ in Equation~\eqref{tor:lattice}.
Lemma~\ref{tor:symmetric-lattice} shows that its degree-two
cohomology is detected on cyclic subgroups.
Lemma~\ref{tor:cyclic-detection} now gives the Hasse principle and
weak approximation for $T$-torsors.

Finally, the solution variety of Equation~\eqref{tor:two-norms}
is a $T$-torsor whenever it is nonempty geometrically: the quotient
of any two solutions has both norms equal to $1$, which is exactly
the defining condition for $T$.
The compatibility of $h$ and $d$ guarantees geometric nonemptiness.
Indeed, over a separable closure write the $2m$ coordinates of $a$
as $a_i^+,a_i^-$, the coordinates of $h$ as $h_i$, and those of
$d$ as $d^+,d^-$.  The equations become
\[
 a_i^+a_i^-=h_i\quad(1\leq i\leq m),\qquad
 \prod_i a_i^+=d^+,\qquad \prod_i a_i^-=d^-.
\]
Choose the $a_i^+$ with product $d^+$ and set
$a_i^-=h_i/a_i^+$.  The final equation follows from
$\prod_i h_i=d^+d^-$.
The asserted local-to-global and approximation statement thus
follows from the result for torsors.
\end{proof}
